% 17_body.tex — Day4 산출물 2/5 (⑰) 운영 이관 체크리스트 · SLA/OLA/UC · 조기 운영 지원(ELS) (빈 템플릿 / 완성 예시 공용 본문)
% 드라이버: 17_운영이관_SLA_ELS_빈템플릿.tex (\wbblanktrue) / 17_운영이관_SLA_ELS_완성예시.tex (\wbblankfalse)
% 내용 출처: PM트랙/Day4_워크북.md §2.3(항목 가이드) §2.4(빈 템플릿) §2.5(온담 P1 완성 예시본)
\documentclass[10pt,a4paper]{article}
\usepackage[a4paper,margin=16mm,top=14mm,bottom=20mm]{geometry}
\def\DocNum{2}
\def\DocTitle{운영 이관 체크리스트 · SLA/OLA/UC · ELS}
\def\DocSub{⑰ Transition Checklist, SLA/OLA/UC \& Early Life Support}
\def\DocVariant{\B{빈 템플릿}{온담 P1 완성 예시}}
\usepackage{wbstyle}
\begin{document}

\docheader
 {\Bg{[정식 명칭과 코드]}{차세대 주문·재고 통합 플랫폼 구축 (ONE ONDAM P1)}}
 {\Bg{[v 계획 / v 실적 — 부속 관계]}{v1.0 2027-02-05 계획 → v1.3 05-21 실적 (SAC 판정 부속)}}
 {\Bg{[인계자 실명]}{P1 PM · 이재현 (인계)\rd{7mm}}}
 {\Bg{[총괄 인수자 실명]}{박준영 서비스운영팀장 (총괄 인수)\rd{7mm}}}

% ------------------------------------------------------------------ 1
\secbar{1. 문서 정보 — 이관 범위 · 인수자 · 기간 · 이관 후 거버넌스}
\guide{이관 대상(시스템·데이터·인터페이스·소스·계약·문서), 인수자(항목마다 다를 수 있으나 \textbf{총괄 인수자는 한 사람}), 인계자, 이관 기간(ELS 시작 ~ G4 — 기간이 없으면 ELS가 끝나지 않는다), 이관 후 거버넌스(CAB·서비스 리뷰 — 프로젝트 CCB와 구분), 입력 문서. 헌장 완료 기준과 SAC 항목이 이관 범위의 뼈대다. 인수자가 “운영본부”이면 인수가 아니다.}
{\footnotesize
\begin{tabularx}{\linewidth}{|L{28mm}|Y|}\hline
\hd{항목} & \hd{내용} \\ \hline
\B{\rd{14mm}}{}이관 대상 & \Bg{[시스템 n종 / 전환 데이터 n TB / 인터페이스 n본 / 소스·형상(하도급 포함) / 계약 n건 / 운영 문서 n종]}{신규 POS·매장 운영 / 주문 허브 / 통합 재고 / 가맹 발주 포털 / 통합 계층(API GW·MDM·SAP 애드온)·전환 데이터 5.1TB / 인터페이스 47본 / 소스·형상(하도급 220 FP 포함) / 계약 5건 이관 / 운영 문서 12종} \\ \hline
\B{\rd{12mm}}{}인수자 & \Bg{[총괄 실명(판정권) / 기술 / 계약 / 보안·개인정보 / 편익 원장 / 인계자]}{총괄 \textbf{박준영}(서비스운영팀장, SAC 판정권) / 기술 박세연 / 계약 강현도·재경팀장 / 보안·개인정보 최영주 / 편익 원장 재무본부(⑳) / 인계자 P1 PM·이재현} \\ \hline
\B{\rd{10mm}}{}기간 & \Bg{[ELS 시작 ~ 종료(n주) → 운영 이관 활동 ~ G4 → 하자보수 기간]}{ELS 2027-03-01 ~ 04-23(8주) → 운영 이관(A20) 04-26 ~ 05-28 → G4. 하자보수 2027-05-29 ~ 2028-05-28(⑱)} \\ \hline
\B{\rd{12mm}}{}이관 후 거버넌스 & \Bg{[변경: CAB 요일·위원장·위원 / 서비스 리뷰 주기·참석 / 프로젝트 CCB 해산 일자]}{변경: 수요일 CAB(위원장 박준영, 위원 인프라·보안·개발운영·수행사 유지보수, 현업 대표 매장운영팀장) / 서비스 리뷰 월 1회(박세연·오정근·한동석) / 프로젝트 CCB 05-28 해산} \\ \hline
\B{\rd{12mm}}{}입력 & \Bg{[런북 실행 기록, SAC, 이슈 로그(미결 n), 감리 시정 확인, 하도급 승인, 계약 변경 로그, 운영 문서 초안, DORA 실측]}{런북 실행 기록(⑯), SAC v1.0, 이슈 로그 v1.0(미결 6), 감리 시정 확인(11-06·G4 05-14), 하도급 승인 기록, 계약 변경 로그, P5 운영 문서 초안, DORA 실측(2026-09)} \\ \hline
\end{tabularx}}

% ------------------------------------------------------------------ 2 (가로) 운영 문서 12종
\landscapeon
\secbar{2. 운영 문서 12종 체크 — 있음 / 없음 / 조건부 \B{}{(서명 05-14, 박준영·운영팀 리더)}}
\guide{12종 각각의 명칭·작성 주체(P1 / P5 — 헌장 §5의 분담을 따르되 목록은 하나로)·버전·상태·\textbf{운영팀 검토 서명(실명·일자)}·미완 시 기한과 책임자. 문서의 존재가 아니라 운영팀이 읽고 서명했는가가 상태다. 프로젝트 중 추가된 절차(아카이브·폴백 배치 파라미터)가 목록에 빠지지 않았는지 본다.}
{\footnotesize\setlength{\tabcolsep}{2.5pt}
\begin{tabularx}{\linewidth}{|C{6mm}|Y|L{20mm}|C{12mm}|L{44mm}|L{34mm}|L{58mm}|}\hline
\hd{\#} & \hd{문서} & \hd{작성 주체} & \hd{버전} & \hd{상태 (있음 / 없음 / 조건부)} & \hd{운영팀 검토 서명} & \hd{비고 · 미완 시 기한 · 책임자} \\ \hline
\ifwbblank
\rd{7mm}1 & \g{[시스템 구성도·아키텍처 정의서]} & \g{[P1]} & & & \g{[실명 일자]} & \\ \hline
\rd{7mm}2 & \g{[운영 매뉴얼(일·주·월 작업)]} & \g{[P5]} & & & & \\ \hline
\rd{7mm}3 & \g{[장애 대응 절차서(운영 런북)]} & & & & & \\ \hline
\rd{7mm}4 & \g{[배치 작업 명세 — 파라미터 포함]} & & & & & \\ \hline
\rd{7mm}5 & \g{[인터페이스 운영 절차(n본, 오류 처리·재전송)]} & & & & & \\ \hline
\rd{7mm}6 & \g{[백업·복구 절차(RPO·RTO·DR)]} & & & & & \\ \hline
\rd{7mm}7 & \g{[모니터링·임계값 정의]} & & & & & \\ \hline
\rd{7mm}8 & \g{[보안 운영 절차(계정·권한·파기)]} & & & & & \\ \hline
\rd{7mm}9 & \g{[배포·형상 절차(빌드·CAB 변경 모델)]} & & & & & \\ \hline
\rd{7mm}10 & \g{[아카이브 접근 절차 등 프로젝트 중 추가된 절차]} & & & \g{[조건부 — 승인 대기]} & & \g{[기한 · 책임자 → 이관 목록]} \\ \hline
\rd{7mm}11 & \g{[서비스 카탈로그·SLA/OLA/UC 서명본]} & & & & & \\ \hline
\rd{7mm}12 & \g{[사용자 매뉴얼·FAQ]} & & & & & \\ \hline
\else
1 & 시스템 구성도·아키텍처 정의서 & P1(수행) & v1.2 & 있음 & 박준영 04-30 & 2차 컷오버 후 센터 원장 반영 \\ \hline
2 & 운영 매뉴얼(일·주·월 작업) & P5 & v1.0 & 있음 & 박준영 05-07 & \\ \hline
3 & 장애 대응 절차서(운영 런북) & P1+P5 & v1.1 & 있음 & 서비스데스크 리더 05-07 & ELS P1-01·02 사례 추가 \\ \hline
4 & 배치 작업 명세(3PL 폴백 21:00·03:00, \textbf{22:30 파라미터}, 일마감, 파기 배치) & P1(수행) & v1.3 & 있음 & 박준영 05-07 & R-16·KE-02 반영 \\ \hline
5 & 인터페이스 운영 절차(47본, 오류 처리·재전송) & P1(수행) & v1.1 & 있음 & 박세연 04-30 & AF-1-02 9본 보완(10-30) 포함 \\ \hline
6 & 백업·복구 절차(RPO 15분·RTO 4h, DR 전환) & P1+인프라 & v1.0 & 있음 & 박준영 04-14 & DR 시험 04-10 실측 3h 20m \\ \hline
7 & 모니터링·임계값 정의(APM 24개) & P1+P5 & v1.2 & 있음 & 박준영 04-30 & ELS 중 12 → 24 \\ \hline
8 & 보안 운영 절차(계정·권한 매트릭스, 파기 주기, 이벤트 중복 검출) & P1+보안 & v1.1 & 있음 & 최영주 05-07 & 03-05 사건 후 갱신 \\ \hline
9 & 배포·형상 절차(빌드·CAB 변경 모델·표준 변경 목록) & P1 PMO & v1.0 & 있음 & 박준영 04-30 & 클린 빌드 재현 05-12 \\ \hline
10 & \textbf{아카이브 접근 절차}(5년 초과 이력, CR-11 조건 ④) & P1+재경 & v1.0 & \textbf{조건부} — 문서·운영팀 서명 완료, 재경팀 승인 대기 & 박준영 05-07 & 승인 기한 \textbf{06-15}, 책임자 재경팀장 → ⑱ 이관 1번 \\ \hline
11 & 서비스 카탈로그·SLA/OLA/UC 서명본 & P5 & v1.0 & 있음 & 박준영·박세연 05-14 & 항목 4 \\ \hline
12 & 사용자 매뉴얼·FAQ(서비스데스크용, 매장용) & P3+P5 & v1.1 & 있음 & 서비스데스크 리더 05-07 & ELS 콜 상위 20건 FAQ 반영 \\ \hline
\fi
\end{tabularx}}

% ------------------------------------------------------------------ 3 Error budget
\secbar[80mm]{3. Error budget — SLA 99.5\% → 월 허용 다운 → 릴리스 속도 환전 \B{}{(day4\_closing.py §1)}}
\guide{가용성 목표 → 월 허용 다운(분, 평균 달 30.4375일) → 계획 정지 처리 규칙 → 비계획(인시던트) 유보와 릴리스 몫의 배분 → 릴리스당 기대 다운(변경 실패율 × 평균 복구 시간) → 월 허용 릴리스 횟수 → 예산 소진 시 규칙(error budget policy). 실패율·복구 시간은 DORA 실측에서 출발해 ELS 목표를 적는다. SLA를 “99.5\%”로만 적고 분으로 환산하지 않으면 CAB는 “안정 vs 기능”의 감정 토론이 된다.}
{\footnotesize
\begin{tabularx}{\linewidth}{|L{50mm}|L{78mm}|Y|}\hline
\hd{항목} & \hd{값} & \hd{산식 · 근거} \\ \hline
\B{\rd{6mm}}{}가용성 목표 & \Bg{[n\% / 월]}{99.5\% / 월 (O-18)} & \Bg{[출처 — 서비스 수준 초안 → SLA 서명]}{P5 서비스 수준 초안 → SLA 서명} \\ \hline
\B{\rd{6mm}}{}월 허용 다운 & \Bg{[n분 = n시간 n분]}{\textbf{219.2분 = 3시간 39분}} & \Bg{[30.4375 × 24 × 60 × (1 $-$ 목표)]}{30.4375일 × 24 × 60 × 0.5\% (28일 달 202분 / 30일 216분 / 31일 223분)} \\ \hline
\B{\rd{6mm}}{}계획 정지 & \Bg{[SLA 산정 제외 / 포함 — 조건]}{SLA 산정 제외} & \Bg{[유지보수 창·컷오버 창 — CAB 승인 건만]}{2차 컷오버 8h(03-27), 월 1회 유지보수 창(일 02:00~04:00) — CAB 승인 건만} \\ \hline
\B{\rd{6mm}}{}비계획(인시던트) 유보 & \Bg{[n\% = n분]}{60\% = 131.5분} & \Bg{[실측 대비]}{ELS 실측 3월 57분 = 유보의 43\%} \\ \hline
\B{\rd{6mm}}{}릴리스 몫 & \Bg{[n\% = n분]}{40\% = \textbf{87.7분}} & \\ \hline
\B{\rd{8mm}}{}릴리스당 기대 다운 (실측) & \Bg{[실패율 n\% × 복구 n분 = n분 → 월 n회(주 n회)]}{12\% × 240분 = \textbf{28.8분} → 월 \textbf{3.0회}(주 0.7회)} & \Bg{[DORA 실측 시점 — 이 상태의 의미]}{DORA 실측 2026-09 — 이 상태로는 주 1회 릴리스 불가} \\ \hline
\B{\rd{8mm}}{}릴리스당 기대 다운 (목표) & \Bg{[n\% × n분 = n분 → 월 n회]}{10\% × 45분 = \textbf{4.5분} → 월 \textbf{19.5회}(주 4.5회)} & \Bg{[복구 시간 단축 수단]}{복구 45분 = KEDB 우회 + 회귀 자동화 + 롤백 스크립트} \\ \hline
\B{\rd{6mm}}{}중간 목표 & \Bg{[하자보수기]}{10\% × 90분 = 9.0분 → 월 9.7회(주 2.2회)} & \Bg{[CAB 릴리스 주기]}{CAB 표준 릴리스 주 1회 + 핫픽스} \\ \hline
\B{\rd{12mm}}{}소진 시 규칙 (error budget policy) & \Bg{[동결 조건 · 해제 조건 · 결정자 · 이의 경로]}{월 누적 다운이 유보 131.5분을 넘으면 표준 릴리스 동결(핫픽스만), 다음 달 초기화 / 릴리스 몫 87.7분 소진 시 CAB 승인 릴리스만} & \Bg{[결정자]}{결정자 박준영, 이의 시 서비스 리뷰} \\ \hline
\B{\rd{8mm}}{}ELS 실적 & \Bg{[월별 다운 · 가용성 · 예산 소진율]}{3월 다운 57분(가용성 99.872\%, 예산 소진 26\%) / 4월 0분(100.000\%) / 8주 합산 99.927\%} & \Bg{[다운 산입 규칙]}{P1-01만 다운 산입; P1-02는 데이터 손상(다운 0)} \\ \hline
\end{tabularx}}

% ------------------------------------------------------------------ 4 SLA/OLA/UC 3층
\clearpage
\secbar{4. SLA / OLA / UC 3층 매핑표 \B{}{(서명 05-14 — SLA 1 · OLA 3 · UC 5)}}
\guide{행 = SLA 항목(가용성·P1 복구·P2 복구·배치 완료·발주 마감 처리·출고 확정·데이터 정정 등), 열 = SLA 목표 / OLA(사내 팀·목표) / UC(외부 계약·목표) / 합산(\textbf{가장 긴 사슬}) / 판정. UC가 없는 행은 “구조적 불가”로 판정하고 SLA 문안을 강등(최선 노력)하거나 UC를 만든다. OLA를 사내 팀에 구두로 부탁하지 않는다. P1 예산 밖 계약(3PL·앱 벤더)을 UC에서 빼지 않는다 — \textbf{한 층이 비면 SLA는 협상이 아니라 산술로 불가능하다.}}
{\footnotesize\setlength{\tabcolsep}{2.5pt}
\begin{tabularx}{\linewidth}{|L{24mm}|L{40mm}|L{44mm}|Y|L{46mm}|L{30mm}|}\hline
\hd{SLA 항목} & \hd{SLA 목표 \B{(고객 ↔ 제공자)}{(현업 오정근·한동석 ↔ IT본부 박세연·박준영)}} & \hd{OLA \B{(사내 팀 · 목표)}{(서비스운영팀 ↔ 인프라팀 / 보안팀 / 개발운영파트)}} & \hd{UC (외부 계약 · 목표)} & \hd{합산 (가장 긴 사슬)} & \hd{판정 (성립 / 재협상 / 강등)} \\ \hline
\ifwbblank
\rd{11mm}\g{① 가용성} & \g{[n\%/월, 계획 정지 처리]} & \g{[착수 n분]} & \g{[CSP n\% / 벤더 패치 SLA / 유지보수]} & \g{[복구 사슬 합 ≤ 허용 다운이면 성립]} & \\ \hline
\rd{11mm}\g{② P1 인시던트 복구} & \g{[응답 n분 / 복구 n시간]} & \g{[접수 → 착수 n분 → 원인 분리 n분]} & \g{[유지보수 착수 n · 복구 n]} & \g{[n + n + n + n = n분 vs 목표]} & \g{[불가 → UC 재협상]} \\ \hline
\rd{11mm}\g{③ P2 인시던트 복구} & & & & & \\ \hline
\rd{11mm}\g{④ 일마감·배치 완료} & & & \g{[3PL 등 UC 유무]} & & \\ \hline
\rd{11mm}\g{⑤ 마감 처리} & & & \g{[— (내부)]} & & \\ \hline
\rd{11mm}\g{⑥ 출고 확정 등 외부 조직 작업} & & & \g{[UC 없음 → 구조적 불가]} & & \g{[최선 노력으로 강등 → 이관]} \\ \hline
\rd{11mm}\g{⑦ 외부 연동 장애 복구} & & & & & \\ \hline
\else
① 가용성 & 99.5\%/월(계획 정지 제외) & 인프라: 인프라 장애 대응 착수 30분 / 개발운영: 앱 장애 착수 30분 & CSP 리전 가용성 99.9\%(CL-01) / API GW·MDM 벤더 패치 SLA 영업일 1일 / 수행사 유지보수 & 다운 원인별 복구 사슬(②) 합이 219분 이내면 성립 & \textbf{성립} (조건: ② 재협상) \\ \hline
② P1 인시던트 복구 & 응답 15분 / \textbf{복구 2시간} & 데스크 접수 → 인프라·개발운영 착수 30분 → 원인 분리 30분 & \textbf{수행사 유지보수 초안: 착수 1시간·복구 4시간} → 재협상 \textbf{착수 30분·복구 90분}(24/7 대기, 연 +0.6억 [추가 설정]) & 초안: 15 + 30 + 30 + 240 = \textbf{315분 > 120} / 재협상 후: 15 + 30 + 30 + 90 = 165분 → 병렬 착수로 \textbf{≤ 120} & 초안 \textbf{불가} → UC 재협상 후 성립 \\ \hline
③ P2 인시던트 복구 & 응답 30분 / 복구 8시간 & 착수 1시간 & 수행사 착수 1시간·복구 4시간 & 30 + 60 + 240 = 330분 ≤ 480 & 성립 \\ \hline
④ 일마감·폴백 배치 완료 & 일마감 23:30, 폴백 배치 1회차 21:40·2회차 03:40 종료 & 개발운영: 배치 실패 재실행 20분 내 & 3PL 대성로지스 \textbf{UC 부재}(물류본부 계약, 준실시간 거부·2027-06 재협상) → “배치 파일 수신 확인 30분”으로 한정 UC 서명(한동석 05-10) & 배치 종료는 P1 측 통제 가능, 수신 이후는 UC 없음 & 배치 종료 항목만 성립 \\ \hline
⑤ 가맹 발주 마감 처리 & 20:00 마감분 21:00까지 출고 지시 생성(CR-09 20:00 운용) & 개발운영 30분 & — (P1 내부) & 성립 & 성립 \\ \hline
⑥ 출고 확정 D+1 06:00 & 발주 → 출고 확정 D+1 06:00 (B-08 12h의 운영 전제) & — & \textbf{3PL UC 없음} (출고 확정은 대성 WMS 소관) & \textbf{구조적 불가} — IT본부가 통제하지 못하는 조직의 작업을 현업에 약속할 수 없음 & \textbf{“최선 노력”으로 강등}, 2027-06 3PL 재협상 후 재상정 → ⑱ 이관 4번 \\ \hline
⑦ 앱 주문 연동 장애 복구 & 4시간 & 개발운영 30분 & 앱 벤더(앱-01, R-18 단가 협상 결과 월 3,600만 원 [추가 설정]) 복구 3시간 & 30 + 180 = 210 ≤ 240 & 성립 \\ \hline
\fi
\end{tabularx}}
\B{\small\g{증명 — 어느 행이 어느 층의 부재로 불가능한가, 언제 발견되어 언제 해소되었는가: }\hrulefill\par}{\small \textbf{증명.} ② 초안 상태에서 SLA P1 복구 2시간은 UC 4시간 하나만으로 이미 불가능했고, ⑥은 UC 자체가 없어 어떤 OLA로도 성립하지 않는다. 한 층이 비면 SLA는 협상이 아니라 산술로 불가능하다 — 그래서 SAC 항목 7은 “SLA 서명”이 아니라 “\textbf{3층 전부 서명}”이다. 재협상(②)과 강등(⑥)은 G3 사전 점검(01-28)에서 발견되어 05-14 전에 끝났다. G4에서 처음 발견되었다면 SLA는 발효되지 못한 채 이관되었을 것이다.\par}

% ------------------------------------------------------------------ 5 ELS 8주
\secbar[80mm]{5. ELS 8주 계획 · 실적 \B{}{(2027-03-01 ~ 04-23, 판정 박준영)}}
\guide{주별 계획(수행사 상주 인력, 워룸, 보고 주기)과 실적(인시던트 P1/P2/P3, 매장 오류율 — \textbf{기준값이 있어야 “좋아졌다”를 말할 수 있다}, 서비스데스크 콜, 핫픽스). 종료 기준은 “2주 연속”처럼 지속 조건으로. 상주 인력 축소는 기준 충족과 연동한다. 2차 컷오버가 ELS 중간에 있으면 그 주의 지표 재상승을 미리 허용 조건으로 둔다.}
{\footnotesize\setlength{\tabcolsep}{2.5pt}
\begin{tabularx}{\linewidth}{|C{8mm}|L{20mm}|L{23mm}|C{9mm}|C{9mm}|C{9mm}|L{30mm}|L{20mm}|C{11mm}|Y|}\hline
\hd{주} & \hd{기간} & \hd{상주 계획 / 실} & \hd{P1} & \hd{P2} & \hd{P3} & \hd{매장 주문 오류율 \B{(기준 → 목표)}{(O-13 1.8\% → 0.5\%)}} & \hd{데스크 콜 \B{(기준/주)}{(기준 400/주)}} & \hd{핫픽스} & \hd{비고} \\ \hline
\ifwbblank
\rd{7mm}1 & & \g{[n / n]} & & & & & & & \g{[주요 인시던트]} \\ \hline
\rd{7mm}2 & & & & & & & & & \g{[1차 하이퍼케어 판정 → 2차 go 조건]} \\ \hline
\rd{7mm}3 & & & & & & & & & \\ \hline
\rd{7mm}4 & & & & & & & & & \\ \hline
\rd{7mm}5 & & & & & & & & & \g{[2차 컷오버 후 재상승 허용]} \\ \hline
\rd{7mm}6 & & & & & & & & & \\ \hline
\rd{7mm}7 & & & & & & & & & \g{[종료 기준 1주차 충족]} \\ \hline
\rd{7mm}8 & & & & & & & & & \g{[2주 연속 → 종료 판정]} \\ \hline
\rd{7mm}합 & & & & & & & & & \\ \hline
\else
1 & 03-01~05 & 20/20 & \textbf{2} & 3 & 9 & 2.4\% & 1,840 & 3 & P1-01(03-03 재고 API 57분), \textbf{P1-02(03-05 포인트 이중차감)}, 정식 라이선스 03-05 \\ \hline
2 & 03-08~12 & 20/20 & 0 & 2 & 8 & 1.3\% & 1,120 & 2 & KE-03 핫픽스 03-12, 1차 하이퍼케어 판정(심각도 2 이상 4건 ≤ 5 → 2차 go 조건 충족) \\ \hline
3 & 03-15~19 & 15/15 & 0 & 1 & 6 & 0.9\% & 780 & 1 & OD-POS 종료 세션 03-19, CAB 첫 회 03-17 \\ \hline
4 & 03-22~26 & 15/15 & 0 & 0 & 5 & 0.7\% & 610 & 1 & 2차 컷오버 준비, 게이트 1 조건 확인 \\ \hline
5 & 03-29~04-02 & 15/18 (2차 후 +3) & 0 & \textbf{2} & 6 & 0.8\% & 690 & 2 & 2차 컷오버 03-27~28, P2-07(03-29 폴백 배치 실패 84분 → KE-02), 지표 재상승 허용 조건 내 \\ \hline
6 & 04-05~09 & 10/10 & 0 & 1 & 4 & 0.5\% & 540 & 1 & 단기 성과 공지 04-05(리드타임 13.6h), DR 전환 시험 04-10 \\ \hline
7 & 04-12~16 & 10/10 & 0 & 0 & 2 & 0.4\% & 470 & 1 & 종료 기준 ①~⑤ 1주차 충족 \\ \hline
8 & 04-19~23 & 6/6 & 0 & 0 & 1 & 0.4\% & 450 & 0 & \textbf{2주 연속 충족 → 04-23 ELS 종료 판정}, 상주 6 → 4(하자보수 체제) \\ \hline
합 & & & \textbf{2} & \textbf{9} & \textbf{41} & & & \textbf{11} & P4 88건 별도(문의·조작 안내) \\ \hline
\fi
\end{tabularx}}
\B{\small\g{종료 기준 ① [2주 연속 P1 0] ② [P2 주 n건 이하] ③ [오류율 n\% 이하] ④ [콜 n/주 이하] ⑤ [심각도 2 backlog 0] / 판정(일자·판정자) / 재상승 허용 조건: }\hrulefill\par}{\small \textbf{종료 기준}: ① 2주 연속 P1 0 ✓(3~8주) ② P2 주 2건 이하 ✓ ③ 매장 주문 오류율 0.5\% 이하 ✓(7·8주 0.4\%) ④ 콜 480/주 이하 ✓(470·450) ⑤ 심각도 2 backlog 0 ✓. 5주차 재상승(P2 2·오류율 0.8\%)은 2차 컷오버 후 허용 조건(“2차 후 1주는 기준 적용 유예”) 내. 판정 박준영 04-23.\par}

% ------------------------------------------------------------------ 6 Incident vs Problem · KEDB
\secbar[80mm]{6. Incident vs Problem · Known Error DB \B{}{(ELS 8주, 인시던트 140 / 문제 9)}}
\guide{인시던트는 “복구되었는가”로 닫고, 문제는 “원인이 제거되었는가”로 닫는다. 같은 원인의 인시던트가 반복되면 문제 1건에 인시던트 n건을 연결한다. Known Error = 원인을 알지만 아직 제거하지 않은(또는 제거하지 않기로 한) 것 — \textbf{우회 방법이 반드시 있어야 한다.} 하자보수(수행사)와 운영 개선(P5)의 경계를 소유자 열로 정한다.}
{\footnotesize
\begin{tabularx}{\linewidth}{|L{30mm}|L{80mm}|Y|}\hline
\hd{구분} & \hd{건수} & \hd{처리} \\ \hline
\B{\rd{7mm}}{}인시던트 & \Bg{[P1 n · P2 n · P3 n · P4 n = n]}{P1 2 · P2 9 · P3 41 · P4 88 = \textbf{140}} & \Bg{[전부 복구 종결 — 평균 복구 시간]}{전부 복구 종결(P1 평균 복구 57분·P2 평균 3h 10m)} \\ \hline
\B{\rd{7mm}}{}문제 레코드 & \Bg{[n = 근본 원인 제거 n + Known Error n]}{9 = 근본 원인 제거 3(재고 API 커넥션 풀·PG 타임아웃 재시도·포털 세션) + \textbf{Known Error 6}} & \Bg{[종결 규칙]}{문제는 원인 제거 시에만 종결} \\ \hline
\end{tabularx}}
\par\vspace{1.5mm}\noindent%
{\footnotesize\setlength{\tabcolsep}{2.5pt}
\begin{tabularx}{\linewidth}{|C{13mm}|L{50mm}|L{36mm}|L{40mm}|Y|L{24mm}|L{24mm}|}\hline
\hd{KE} & \hd{증상} & \hd{우회} & \hd{근본 원인} & \hd{상태 \B{(일자)}{(05-21)}} & \hd{소유자} & \hd{연결 인시던트} \\ \hline
\ifwbblank
\rd{9mm}KE-01 & & \g{[반드시 있어야 한다]} & & \g{[유지 / 근본 제거 일자 / 보류 사유]} & \g{[실명]} & \g{[P1-nn, P3 n건]} \\ \hline
\rd{9mm}KE-02 & & & & & & \\ \hline
\rd{9mm}KE-03 & & & & & & \\ \hline
\rd{9mm}KE-04 & & & & & & \\ \hline
\rd{9mm}KE-05 & & & & & & \\ \hline
\rd{9mm}KE-06 & & & & & & \\ \hline
\else
KE-01 & 구형 단말(2013년 이전 기종 35점 중 잔여 9점) 영수증 인쇄 지연 8초 & 재부팅, 인쇄 큐 초기화 & 펌웨어 드라이버 호환(CR-14 패치 범위 밖) & Known Error 유지 — 단말 교체 협의(⑱ 이관 2번) & 오정근 & P3 7건 \\ \hline
KE-02 & 폴백 배치 1회차 40분 초과 시 2회차와 겹쳐 재고 정합 오류(R-16) & 1회차 시작을 22:30로(파라미터) & 배치 대상 증가(2차 후 센터 재고 포함) & 우회 적용 04-01, 근본 해소는 3PL 준실시간 전환 시(2027-06 재협상) & 한동석·박준영 & P2-07, P3 3건 \\ \hline
KE-03 & 반품 처리 시 멤버십 취소 이벤트 중복 → 포인트 이중차감(D-0902 계열) & 중복 검출 배치(일 1회) + 자동 복구 & 이벤트 멱등 키 누락 & \textbf{근본 제거 03-12 핫픽스}, 재발 감시 배치 유지로 KEDB 잔류 & 수행사 유지보수 & \textbf{P1-02}, P3 2건 \\ \hline
KE-04 & SAP 출고 확정 IF-SAP-12 타임아웃 시 재시도 3회 후 실패(D-0944) & 수동 재전송 화면 & SAP ECC 야간 배치와 경합 & Known Error — ECC 교체(2027-12 EoM) 전까지 유지 & 박세연 & P3 9건 \\ \hline
KE-05 & 엑셀 발주 구형 xls 인코딩 수량 0(D-0931) & xlsx 변환 안내, 업로드 전 검증 메시지 & 구형 라이브러리 & Known Error — 사용률 1\% 미만, 제거 보류(CAB 04-14) & 박준영 & P4 12건 \\ \hline
KE-06 & 옴니 매장 간 재고 이동 요청 시 재고 예약 반영 15분 지연(폴백 창) & 안내 문구·재고 재조회 & 3PL 폴백 배치 주기 & Known Error — KE-02와 동일 원인 계열 & 오정근 & P3 4건 \\ \hline
\fi
\end{tabularx}}

% ------------------------------------------------------------------ 7 개인정보 사건 판단 기록
\secbar[100mm]{7. 개인정보 사건 판단 기록 — 조문 단정 · 적용 비단정 \B{}{— P1-02 가맹 포인트 이중차감}}
\guide{사실(인지 시각 — \textbf{72시간의 기산점을 분 단위로}, 규모, 영향 정보주체 수), 적용 조문(정확히 인용 — 단정형), 쟁점(이 사건이 조문의 어느 요건에 해당하는가 — 단정하지 않음), 검토 의견(실명·일시, 단정하지 않는 부분 명시), 결정(통지와 신고를 구분, 각각의 시각), 근거·기록, 교훈. “유출 아님”으로 자체 판정하고 기록을 남기지 않는 것과, 조문과 해석을 섞어 “법상 신고 대상”이라 단정하는 것이 양쪽의 오류다.}
{\footnotesize
\begin{tabularx}{\linewidth}{|L{36mm}|Y|}\hline
\hd{항목} & \hd{내용} \\ \hline
\B{\rd{12mm}}{}사실 & \Bg{[YYYY-MM-DD hh:mm 인지(경위) / 기간·건수·영향 정보주체 n명(집단별) / 원인(KE) / 외부 접근·유출 정황 확인(로그·시각)]}{2027-03-05(금) \textbf{14:20 인지}(서비스데스크 동일 유형 문의 31건 집계 → P1 선언). 03-01~05 반품 처리 3,412건에서 멤버십 포인트 취소가 2회 적용, 영향 회원 \textbf{2,987명}(가맹 매장 회원 2,140·직영 847), 오차감 합계 1,284만 포인트. 원인 KE-03(이벤트 중복). 외부 접근·유출 정황 없음(보안팀 로그 03-05 20:00 확인)} \\ \hline
\B{\rd{12mm}}{}적용 조문 (단정) & \Bg{[법률 번호·시행일 — ‘유출등’ 정의(위조·변조·훼손 포함), 인지 후 72시간 이내 통지·신고, 신고 요건은 시행령·고시]}{개인정보 보호법(법률 제21445호, 2026-03-10 공포·\textbf{2026-09-11 시행}): 통지·신고 대상인 ‘유출등’의 정의에 분실·도난·유출 외 \textbf{위조·변조·훼손}이 포함된다. 유출등을 알게 된 때부터 \textbf{72시간 이내} 정보주체 통지·보호위원회 신고(신고 대상 요건은 시행령·고시)} \\ \hline
\B{\rd{12mm}}{}쟁점 (적용 — 단정하지 않음) & \Bg{[① 이 사건이 ‘훼손’에 해당하는가 — 해석례 유무 ② 규모 요건 ③ 복구 후 신고 의무 유지 여부]}{① 시스템 오류로 회원의 포인트 잔액이 사실과 다르게 기록된 것이 ‘훼손’에 해당하는가 — 개정법 시행 후 해석례·결정례 부재 ② 해당한다면 영향 2,987명이 신고 대상 규모 요건에 해당하는가(시행령 기준 확인 필요) ③ 복구 완료 후에도 신고 의무가 유지되는가} \\ \hline
\B{\rd{16mm}}{}검토 의견 & \Bg{[CPO 실명, 일시, 서면 — 해석 여지가 있는 부분과 보수적 권고, 통지·복구는 규제 해당 여부와 무관하게 즉시]}{\textbf{최영주 CPO, 03-06(토) 11:00 서면}: “①은 해석 여지가 있다. 포인트 잔액은 회원 식별 정보와 결합된 개인정보이고 개정법이 ‘훼손’을 명시한 취지(외부 유출 없이도 무결성 침해를 포착)에 비추어 해당 가능성을 배제할 수 없다. 단정할 수 없으므로 \textbf{보수적으로 72시간 내 신고를 권고}한다. 정보주체 통지·복구는 규제 해당 여부와 무관하게 즉시 실시한다. 과징금 기준이 개정법으로 상향된 점을 고려한다.”} \\ \hline
\B{\rd{14mm}}{}결정 & \Bg{[결정자 실명·일시 / ① 통지·복구 완료 시각 ② 신고 시각(기산 후 n시간 n분) ③ 이해관계자 설명 ④ 재발 방지]}{정미란·최영주·박준영·P1 PM 03-06 14:00. ① 회원 개별 통지(앱 푸시·문자)·포인트 자동 복구 \textbf{03-06 18:00 완료} ② \textbf{개인정보보호위원회 신고 03-07(일) 16:00} — 기산 후 49시간 40분 ③ 협의회(서정필) 서면 설명 03-08 ④ 재발 방지: 이벤트 멱등 키(03-12 핫픽스), 중복 검출 배치, 운영 문서 8종 갱신} \\ \hline
\B{\rd{8mm}}{}근거 · 기록 & \Bg{[티켓 번호·인지 시각, 로그 확인서, 검토 의견서, 신고 접수증, 통지 발송 로그, 사본 수신]}{인지 시각 데스크 티켓 \#ELS-0417 14:20, 보안 로그 확인서, CPO 검토 의견서, 신고 접수증, 통지 발송 로그 2,987건. 사본 윤태호(내부감사)} \\ \hline
\B{\rd{8mm}}{}교훈 & \Bg{[L-nn — 운영 문서에 고정할 시한 절차]}{L-14(⑲): 개인정보 사건 판단 절차를 운영 문서 8종에 “인지 시각 기록 → CPO 24시간 내 의견 → 통지 즉시 → 신고 판단 48시간 내”로 고정. 규제 적용을 단정하지 않되 시한은 단정한다} \\ \hline
\end{tabularx}}

% ------------------------------------------------------------------ 8 사용자 채택 · ADKAR
\secbar[80mm]{8. 사용자 채택 지표 · ADKAR 장벽 진단 \B{}{(P3 공동, 데이터 P3 교육 대장·포털 로그)}}
\guide{집단을 나눈다(직영·가맹·서비스데스크·물류센터). 채택 지표의 기준값·목표·ELS 주별 실적. 장벽은 A(인식)·D(욕구)·K(지식)·A(능력)·R(강화) 중 \textbf{가장 낮은 단계}로 진단하고 그 단계에 맞는 조치만 적는다(지식 부족에 인식 캠페인을 하지 않는다). Kotter의 단기 성과를 지표로 공지한다. Bridges의 “끝냄(ending)”을 의식적으로 처리한다 — 구 시스템을 조용히 끄지 않는다. 채택을 교육 이수율로만 재지 않는다.}
{\footnotesize\setlength{\tabcolsep}{2.5pt}
\begin{tabularx}{\linewidth}{|L{28mm}|L{44mm}|C{11mm}|C{11mm}|C{13mm}|L{56mm}|Y|}\hline
\hd{집단} & \hd{지표 (기준값 → 목표)} & \hd{2주} & \hd{4주} & \hd{8주} & \hd{장벽 진단 (A/D/K/A/R)} & \hd{조치 (P1 / P3)} \\ \hline
\ifwbblank
\rd{12mm}\g{[가맹 점주]} & \g{[전화·팩스 발주 비율 n\% → n\%]} & & & & \g{[가장 낮은 단계 + 근거]} & \g{[그 단계의 조치만]} \\ \hline
\rd{12mm}\g{[직영 매장]} & \g{[신규 절차 준수율]} & & & & & \\ \hline
\rd{12mm}\g{[서비스데스크]} & \g{[1차 해결률]} & & & & & \\ \hline
\rd{12mm}\g{[물류센터]} & \g{[사용률]} & & & & & \\ \hline
\rd{12mm} & & & & & & \\ \hline
\else
가맹 점주 (388점, 331명) & 전화·팩스·엑셀 발주 비율 39\% → 15\% 이하 & 24\% & 16\% & \textbf{11\%} & \textbf{K/A} — 구형 단말 잔여 9점, 60대 이상 점주 27\%가 포털 조작에 시간 소요 & P3: 지역매니저 방문 교육 2회차(04월), 포털 “간편 재발주” 기능(R5 05-07). D(욕구) 장벽은 부속서 동의 91\%로 대부분 해소 \\ \hline
가맹 점주 & 포털 마감 20:00 준수율 → 95\% & 81\% & 90\% & 96\% & R(강화) — 마감 놓치면 다음날 출고 불가 경험이 강화 & 마감 30분 전 알림 \\ \hline
직영 매장 (227점) & 반품 신규 절차(멤버십 취소 확인) 준수율 → 95\% & 78\% & 88\% & 94\% & \textbf{R} — 절차는 알지만 바쁜 시간대 생략 & POS 화면 강제 확인 단계(핫픽스 03-12와 함께), 지역매니저 점검 \\ \hline
서비스데스크 (11명) & 1차 해결률 → 70\% & 52\% & 64\% & 73\% & K — FAQ 부족 & FAQ 20건(문서 12종), KEDB 우회 방법 공유 \\ \hline
물류센터 (246명, 2차 후) & 센터 재고 조회 신 플랫폼 사용률 → 90\% & — & — & 87\%(4주) & A(인식)/D — 노조(김미르) “처리량 목표와 인력” 우려가 시스템 사용 저항으로 & 한동석: 재고 정확도 리포트는 인력 평가에 쓰지 않음을 서면(04-12), P2 편익과 분리 확인 \\ \hline
\fi
\end{tabularx}}
{\footnotesize
\begin{tabularx}{\linewidth}{|L{28mm}|Y|}\hline
\B{\rd{9mm}}{}단기 성과 공지 (Kotter) & \Bg{[일자 · 매체 · 내용 — 숫자로. 목표 판정 시점과 구분]}{04-05 협의회 공문 “가맹 발주 리드타임 38시간 → \textbf{13.6시간}(4월 1주 평균)” — 목표 12시간은 M+12 판정 → ⑳ B-08 첫 측정} \\ \hline
\B{\rd{9mm}}{}끝냄 처리 (Bridges) & \Bg{[구 시스템 종료 세션 일자·주관·형식 — “없어진 시스템”이 아니라 “임무를 마친 시스템”으로]}{03-19 OD-POS 2.1 종료 세션 — 박세연(2011년 원 설계자) 주관, 15년 운영 기록 아카이브, 개발자 초청. 강윤지 후임 지역매니저 “이번엔 없어지지 않는다”의 답} \\ \hline
\end{tabularx}}
\landscapeoff

% ------------------------------------------------------------------ 9 이슈 종결 · 인수 서명
\secbar[80mm]{9. 이슈 종결 확인 · 인수 서명}
\guide{Day3 이슈 로그 미결 건의 최종 처리 — “종결” 또는 “운영 리스크로 이관(⑱ 항목 5·6)” 둘 중 하나. 진행 중인 채로 G4를 넘기지 않는다. 잔여 작업의 소유자 실명. 총괄 인수 서명과 기술·계약·보안 인수 서명의 일자.}
{\footnotesize
\begin{tabularx}{\linewidth}{|L{30mm}|L{22mm}|Y|L{34mm}|}\hline
\hd{IS} & \hd{최종 처리} & \hd{근거 · 일자} & \hd{잔여 소유자} \\ \hline
\ifwbblank
\rd{8mm}\g{[IS-nn 제목]} & \g{[종결 / 운영 이관]} & & \g{[실명 (이관 목록 번호)]} \\ \hline
\rd{8mm} & & & \\ \hline
\rd{8mm} & & & \\ \hline
\rd{8mm} & & & \\ \hline
\rd{8mm} & & & \\ \hline
\rd{8mm} & & & \\ \hline
\else
IS-04 슈퍼유저 이탈 & 종결 & P3 재선정 11-13, 후임 지역매니저 검토회 참여, UAT 참여자 확정 12월 & — \\ \hline
IS-06 가맹 동의 71\% & 종결 & 12-22 동의 91\%(353점), 미동의 35점 CR-14 2차 배포 → 잔여 9점 협의 & 오정근(⑱ 이관 2번) \\ \hline
IS-07 감리 귀속 이의 & 종결 & 정하늘 회신 10-08 “발주사 귀속”으로 정정, 시정 확인 11-06, 변경협의체 10-14 대기 비용 0.18 판정 & — \\ \hline
IS-08 보안 High 2건 & 종결 & 조치 10-16, 재진단 10-23 High 0, 2회차(01-20) High 0 & — \\ \hline
IS-09 22:00 마감 요구 & 종결 & CR-09 파라미터 12-04 릴리스, 운용 시각 CCB 01-21 20:00 확정, 협의회 회신 01-25 & 한동석(⑱ 이관 3번 — 2027-09 CAB 재상정) \\ \hline
IS-10 클라우드 사용량 & \textbf{운영 이관} & 부하시험(01월) 후 사이징 확정 — 유보 3.0 복귀 불요(반납), 운영 사용량 월 0.42 → P5 운영비 반영 & 박준영(P5 예산) \\ \hline
\fi
\end{tabularx}}
\vspace{3mm}
\noindent{\footnotesize
\begin{tabularx}{\linewidth}{|Y|Y|Y|Y|Y|}\hline
\hd{총괄 인수 (서명 · 일자)} & \hd{기술 인수 (서명 · 일자)} & \hd{계약 인수 (서명 · 일자)} & \hd{보안 인수 (서명 · 일자)} & \hd{인계 (서명 · 일자)} \\ \hline
\Bg{[서비스운영팀장]}{박준영 2027-05-21 (SAC 판정과 동시)\rd{10mm}} & \Bg{[IT본부장]}{박세연 05-12} & \Bg{[구매팀장]}{강현도 05-26} & \Bg{[CPO]}{최영주 05-07} & \Bg{[P1 PM · 수행사 PM]}{P1 PM · 이재현 05-21} \\ \hline
\end{tabularx}}
\end{document}
